\documentclass[10pt,a4paper]{amsart}
\usepackage{amsmath,amssymb,amsthm,graphicx}
\usepackage{verbatim}
\newenvironment{reference}{\comment}{\endcomment}
%\newenvironment{reference}{}{}
\newenvironment{slogan}{\comment}{\endcomment}
\newenvironment{history}{\comment}{\endcomment}

% For commutative diagrams we use Xy-pic
\usepackage[all]{xy}

% We use 2cell for 2-commutative diagrams.
\xyoption{2cell}
% Theorem environments.
%
\theoremstyle{plain}
\newtheorem{theorem}[subsection]{Theorem}
\newtheorem{proposition}[subsection]{Proposition}
\newtheorem{lemma}[subsection]{Lemma}

\theoremstyle{definition}
\newtheorem{definition}[subsection]{Definition}
\newtheorem{example}[subsection]{Example}
\newtheorem{exercise}[subsection]{Exercise}
\newtheorem{situation}[subsection]{Situation}

\theoremstyle{remark}
\newtheorem{remark}[subsection]{Remark}
\newtheorem{remarks}[subsection]{Remarks}

\numberwithin{equation}{subsection}

% Macros
%
\def\lim{\mathop{\mathrm{lim}}\nolimits}
\def\colim{\mathop{\mathrm{colim}}\nolimits}
\def\Spec{\mathop{\mathrm{Spec}}}
\def\Hom{\mathop{\mathrm{Hom}}\nolimits}
\def\Ext{\mathop{\mathrm{Ext}}\nolimits}
\def\SheafHom{\mathop{\mathcal{H}\!\mathit{om}}\nolimits}
\def\SheafExt{\mathop{\mathcal{E}\!\mathit{xt}}\nolimits}
\def\Sch{\mathit{Sch}}
\def\Mor{\mathop{\mathrm{Mor}}\nolimits}
\def\Ob{\mathop{\mathrm{Ob}}\nolimits}
\def\Sh{\mathop{\mathit{Sh}}\nolimits}
\def\NL{\mathop{N\!L}\nolimits}
\def\CH{\mathop{\mathrm{CH}}\nolimits}
\def\proetale{{pro\text{-}\acute{e}tale}}
\def\etale{{\acute{e}tale}}
\def\QCoh{\mathit{QCoh}}
\def\Ker{\mathop{\mathrm{Ker}}}
\def\Im{\mathop{\mathrm{Im}}}
\def\Coker{\mathop{\mathrm{Coker}}}
\def\Coim{\mathop{\mathrm{Coim}}}

% Boxtimes
%
\DeclareMathSymbol{\boxtimes}{\mathbin}{AMSa}{"02}

%
% Macros for moduli stacks/spaces
%
\def\QCohstack{\mathcal{QC}\!\mathit{oh}}
\def\Cohstack{\mathcal{C}\!\mathit{oh}}
\def\Spacesstack{\mathcal{S}\!\mathit{paces}}
\def\Quotfunctor{\mathrm{Quot}}
\def\Hilbfunctor{\mathrm{Hilb}}
\def\Curvesstack{\mathcal{C}\!\mathit{urves}}
\def\Polarizedstack{\mathcal{P}\!\mathit{olarized}}
\def\Complexesstack{\mathcal{C}\!\mathit{omplexes}}
% \Pic is the operator that assigns to X its picard group, usage \Pic(X)
% \Picardstack_{X/B} denotes the Picard stack of X over B
% \Picardfunctor_{X/B} denotes the Picard functor of X over B
\def\Pic{\mathop{\mathrm{Pic}}\nolimits}
\def\Picardstack{\mathcal{P}\!\mathit{ic}}
\def\Picardfunctor{\mathrm{Pic}}
\def\Deformationcategory{\mathcal{D}\!\mathit{ef}}
\usepackage{cleveref}
\crefname{originallemma}{original lemma}{original lemmas}
\Crefname{originallemma}{Original lemma}{Original lemmas}
\makeatletter
\newlabel{lemma-compute-symbol-M}{{02Q5}{source}{}{sourcepointer-1}{}}
\newlabel{lemma-compute-symbol-M@cref}{{[originallemma][0][]02Q5}{[1][0][]source}{}{}{}}
\makeatother
\crefname{originallemma}{original lemma}{original lemmas}
\Crefname{originallemma}{Original lemma}{Original lemmas}
\makeatletter
\newlabel{lemma-determinant-quotient-ring}{{02PC}{source}{}{sourcepointer-2}{}}
\newlabel{lemma-determinant-quotient-ring@cref}{{[originallemma][0][]02PC}{[1][0][]source}{}{}{}}
\makeatother
\crefname{originallemma}{original lemma}{original lemmas}
\Crefname{originallemma}{Original lemma}{Original lemmas}
\makeatletter
\newlabel{lemma-symbol-when-one-is-a-unit}{{02Q2}{source}{}{sourcepointer-3}{}}
\newlabel{lemma-symbol-when-one-is-a-unit@cref}{{[originallemma][0][]02Q2}{[1][0][]source}{}{}{}}
\makeatother
\crefname{originallemma}{original lemma}{original lemmas}
\Crefname{originallemma}{Original lemma}{Original lemmas}
\makeatletter
\newlabel{lemma-symbol-when-equal}{{0AY9}{source}{}{sourcepointer-4}{}}
\newlabel{lemma-symbol-when-equal@cref}{{[originallemma][0][]0AY9}{[1][0][]source}{}{}{}}
\makeatother
\crefname{originallemma}{original lemma}{original lemmas}
\Crefname{originallemma}{Original lemma}{Original lemmas}
\makeatletter
\newlabel{lemma-Noetherian-domain-dim-1-two-elements}{{02Q7}{source}{}{sourcepointer-5}{}}
\newlabel{lemma-Noetherian-domain-dim-1-two-elements@cref}{{[originallemma][0][]02Q7}{[1][0][]source}{}{}{}}
\makeatother
\crefname{originallemma}{original lemma}{original lemmas}
\Crefname{originallemma}{Original lemma}{Original lemmas}
\makeatletter
\newlabel{lemma-pre-symbol}{{02PX}{source}{}{sourcepointer-6}{}}
\newlabel{lemma-pre-symbol@cref}{{[originallemma][0][]02PX}{[1][0][]source}{}{}{}}
\makeatother
\crefname{originallemma}{original lemma}{original lemmas}
\Crefname{originallemma}{Original lemma}{Original lemmas}
\makeatletter
\newlabel{lemma-symbol-compare-modules}{{02Q4}{source}{}{sourcepointer-7}{}}
\newlabel{lemma-symbol-compare-modules@cref}{{[originallemma][0][]02Q4}{[1][0][]source}{}{}{}}
\makeatother
\crefname{originallemma}{original lemma}{original lemmas}
\Crefname{originallemma}{Original lemma}{Original lemmas}
\makeatletter
\newlabel{lemma-periodic-determinant-easy-case}{{02PN}{source}{}{sourcepointer-8}{}}
\newlabel{lemma-periodic-determinant-easy-case@cref}{{[originallemma][0][]02PN}{[1][0][]source}{}{}{}}
\makeatother
\crefname{originallemma}{original lemma}{original lemmas}
\Crefname{originallemma}{Original lemma}{Original lemmas}
\makeatletter
\newlabel{lemma-periodic-determinant}{{02PO}{source}{}{sourcepointer-9}{}}
\newlabel{lemma-periodic-determinant@cref}{{[originallemma][0][]02PO}{[1][0][]source}{}{}{}}
\makeatother
\crefname{originallemma}{original lemma}{original lemmas}
\Crefname{originallemma}{Original lemma}{Original lemmas}
\makeatletter
\newlabel{lemma-periodic-determinant-sign}{{02PM}{source}{}{sourcepointer-10}{}}
\newlabel{lemma-periodic-determinant-sign@cref}{{[originallemma][0][]02PM}{[1][0][]source}{}{}{}}
\makeatother
\title{The Steinberg relation for the symbol of a one-dimensional module over a Noetherian local ring}
\author{The Stacks Project Authors}
\date{stacks-project:chow.tex:local-snapshot-20261012; GNU FDL 1.2 or later}
\begin{document}
\maketitle
\noindent\textit{Editorial note.} All delivered mathematical text is the author's own, copied from the pinned Stacks Project chapter file chow.tex; no statement or proof step was rewritten, reordered or removed.\par\medskip
\noindent\textit{Editorial note.} Preamble provenance: the pinned chapter file opens with an input of the shared preamble and carries no preamble text of its own, so the project's single Stacks-wide preamble staged at sources/stacks-chow/preamble.tex is excerpted verbatim as two spans, lines 14-18 (Xy-pic setup) and lines 154-217 (the native theorem declarations and mathematical macros). Nothing was added, removed or redefined; the wrapper class and the amsmath/amssymb/amsthm/graphicx base packages are supplied by the assembler.\par\medskip
\noindent\textit{Editorial note.} The counted statement and proof contain no cross reference, no citation and no included file, so the delivered body is self contained.\par\medskip
\noindent\textit{Editorial note.} Canonical location of the counted proposition: Stacks tag 02Q8 (sources/stacks-chow/pinned-original-source.tex line 17660).\par\medskip
\noindent\textit{Allowed standard prerequisites (editorial).}\begin{itemize}
\item The symbol d\_M(a, b) of a pair of elements on a one-dimensional module and its multiplicativity. Same chapter, Lemma labelled lemma-compute-symbol-M and Lemma labelled lemma-symbol-compare-modules, cited at proof lines 17672-17673 and 17716-17718, with the relative symbols of Lemma labelled lemma-pre-symbol at proof lines 17705-17709. Supplies the object whose relation is asserted, the reduction to prime cyclic modules, and the comparison between relative and absolute symbols.
\item Identification of determinants over A and over a quotient A/q. Same chapter, Lemma labelled lemma-determinant-quotient-ring, cited at proof lines 17683-17685. Justifies replacing A by the one-dimensional domain A/q.
\item Two-element factorizations in a one-dimensional local Noetherian domain. Same chapter, Lemma labelled lemma-Noetherian-domain-dim-1-two-elements, cited at proof lines 17693-17695. Produces the extension A' and the factorizations a = t a', b = t b' with (a', b') = A'.
\item Periodic determinant computations and the sign of d\textasciicircum{}A\_\{A'\}(b', b'). Same chapter, Lemmas labelled lemma-periodic-determinant-easy-case, lemma-periodic-determinant and lemma-periodic-determinant-sign, cited at proof lines 17736-17758. Evaluates the three relative symbols through the displayed decompositions and supplies the final sign.
\end{itemize}\medskip
\ifdefined\hypertarget\hypertarget{sourcepointer-1}{}\fi
\noindent\textit{Source reference.} Original lemma 02Q5: sources/stacks-chow/pinned-original-source.tex, source lines 17577--17577. Same-chapter pointer into the author environment lemma at source line 17577; invoked at proof line 17673 of the counted proof. It is retained only as a navigation pointer; every load-bearing step is inside the delivered span.\par\medskip
\ifdefined\hypertarget\hypertarget{sourcepointer-2}{}\fi
\noindent\textit{Source reference.} Original lemma 02PC: sources/stacks-chow/pinned-original-source.tex, source lines 16434--16434. Same-chapter pointer into the author environment lemma at source line 16434; invoked at proof line 17684 of the counted proof. It is retained only as a navigation pointer; every load-bearing step is inside the delivered span.\par\medskip
\ifdefined\hypertarget\hypertarget{sourcepointer-3}{}\fi
\noindent\textit{Source reference.} Original lemma 02Q2: sources/stacks-chow/pinned-original-source.tex, source lines 17483--17483. Same-chapter pointer into the author environment lemma at source line 17483; invoked at proof line 17689 of the counted proof. It is retained only as a navigation pointer; every load-bearing step is inside the delivered span.\par\medskip
\ifdefined\hypertarget\hypertarget{sourcepointer-4}{}\fi
\noindent\textit{Source reference.} Original lemma 0AY9: sources/stacks-chow/pinned-original-source.tex, source lines 17472--17472. Same-chapter pointer into the author environment lemma at source line 17472; invoked at proof line 17690 of the counted proof. It is retained only as a navigation pointer; every load-bearing step is inside the delivered span.\par\medskip
\ifdefined\hypertarget\hypertarget{sourcepointer-5}{}\fi
\noindent\textit{Source reference.} Original lemma 02Q7: sources/stacks-chow/pinned-original-source.tex, source lines 566--566. Same-chapter pointer into the author environment lemma at source line 566; invoked at proof line 17695 of the counted proof. It is retained only as a navigation pointer; every load-bearing step is inside the delivered span.\par\medskip
\ifdefined\hypertarget\hypertarget{sourcepointer-6}{}\fi
\noindent\textit{Source reference.} Original lemma 02PX: sources/stacks-chow/pinned-original-source.tex, source lines 17371--17371. Same-chapter pointer into the author environment lemma at source line 17371; invoked at proof line 17706 of the counted proof. It is retained only as a navigation pointer; every load-bearing step is inside the delivered span.\par\medskip
\ifdefined\hypertarget\hypertarget{sourcepointer-7}{}\fi
\noindent\textit{Source reference.} Original lemma 02Q4: sources/stacks-chow/pinned-original-source.tex, source lines 17533--17533. Same-chapter pointer into the author environment lemma at source line 17533; invoked at proof line 17710 of the counted proof. It is retained only as a navigation pointer; every load-bearing step is inside the delivered span.\par\medskip
\ifdefined\hypertarget\hypertarget{sourcepointer-8}{}\fi
\noindent\textit{Source reference.} Original lemma 02PN: sources/stacks-chow/pinned-original-source.tex, source lines 16696--16696. Same-chapter pointer into the author environment lemma at source line 16696; invoked at proof line 17731 of the counted proof. It is retained only as a navigation pointer; every load-bearing step is inside the delivered span.\par\medskip
\ifdefined\hypertarget\hypertarget{sourcepointer-9}{}\fi
\noindent\textit{Source reference.} Original lemma 02PO: sources/stacks-chow/pinned-original-source.tex, source lines 16730--16730. Same-chapter pointer into the author environment lemma at source line 16730; invoked at proof line 17732 of the counted proof. It is retained only as a navigation pointer; every load-bearing step is inside the delivered span.\par\medskip
\ifdefined\hypertarget\hypertarget{sourcepointer-10}{}\fi
\noindent\textit{Source reference.} Original lemma 02PM: sources/stacks-chow/pinned-original-source.tex, source lines 16676--16676. Same-chapter pointer into the author environment lemma at source line 16676; invoked at proof line 17755 of the counted proof. It is retained only as a navigation pointer; every load-bearing step is inside the delivered span.\par\medskip
\begin{lemma}
\label{lemma-symbol-is-steinberg-prepare}
Let $A$ be a Noetherian local ring.
Let $a, b \in A$.
Let $M$ be a finite $A$-module of dimension $1$ on
which each of $a$, $b$, $b - a$ are nonzerodivisors.
Then
$$
d_M(a, b - a)d_M(b, b) = d_M(b, b - a)d_M(a, b)
$$
in $\kappa^*$.
\end{lemma}

\begin{proof}
By Lemma \ref{lemma-compute-symbol-M} it suffices to show the relation when
$M = A/\mathfrak q$ for some prime $\mathfrak q \subset A$ with
$\dim(A/\mathfrak q) = 1$.

\medskip\noindent
In case $M = A/\mathfrak q$ we may replace $A$ by $A/\mathfrak q$ and
$a, b$ by their images in $A/\mathfrak q$. Hence we may assume $A = M$
and $A$ a local Noetherian domain of dimension $1$. The reason is
that the residue field $\kappa$ of $A$ and $A/\mathfrak q$ are
the same and that for any $A/\mathfrak q$-module $M$ the determinant
taken over $A$ or over $A/\mathfrak q$ are canonically identified.
See Lemma \ref{lemma-determinant-quotient-ring}.

\medskip\noindent
It suffices to show the relation when both
$a, b$ are in the maximal ideal. Namely, the case where one
or both are units follows from Lemmas \ref{lemma-symbol-when-one-is-a-unit}
and \ref{lemma-symbol-when-equal}.

\medskip\noindent
Choose an extension $A \subset A'$ and factorizations
$a = ta'$, $b = tb'$ as in
Lemma \ref{lemma-Noetherian-domain-dim-1-two-elements}.
Note that also $b - a = t(b' - a')$ and that
$A' = (a', b') = (a', b' - a') = (b' - a', b')$.
Here and in the following we think of $A'$ as an $A$-module and
$a, b, a', b', t$ as $A$-module endomorphisms of $A'$.
We will use the notation $d^A_{A'}(a', b')$ and so on to indicate
$$
d^A_{A'}(a', b')
=
\det\nolimits_\kappa(A'/a'b'A', a', b')
$$
which is defined by Lemma \ref{lemma-pre-symbol}. The upper index ${}^A$
is used to distinguish this from the already defined symbol
$d_{A'}(a', b')$ which is different (for example because it has values
in the residue field of $A'$ which may be different from $\kappa$).
By Lemma \ref{lemma-symbol-compare-modules} we see that
$d_A(a, b) = d^A_{A'}(a, b)$,
and similarly for the other combinations.
Using this and multiplicativity we see that it suffices to prove
$$
d^A_{A'}(a', b' - a') d^A_{A'}(b', b')
=
d^A_{A'}(b', b' - a') d^A_{A'}(a', b')
$$
Now, since $(a', b') = A'$ and so on we have
$$
\begin{matrix}
A'/(a'(b' - a')) & \cong & A'/(a') \oplus A'/(b' - a') \\
A'/(b'(b' - a')) & \cong & A'/(b') \oplus A'/(b' - a') \\
A'/(a'b') & \cong & A'/(a') \oplus A'/(b')
\end{matrix}
$$
Moreover, note that multiplication by $b' - a'$ on
$A/(a')$ is equal to multiplication by $b'$, and that
multiplication by $b' - a'$ on $A/(b')$ is equal to multiplication by $-a'$.
Using Lemmas
\ref{lemma-periodic-determinant-easy-case} and
\ref{lemma-periodic-determinant}
we conclude
$$
\begin{matrix}
d^A_{A'}(a', b' - a') & = &
\det\nolimits_\kappa(b'|_{A'/(a')})^{-1}
\det\nolimits_\kappa(a'|_{A'/(b' - a')}) \\
d^A_{A'}(b', b' - a') & = &
\det\nolimits_\kappa(-a'|_{A'/(b')})^{-1}
\det\nolimits_\kappa(b'|_{A'/(b' - a')}) \\
d^A_{A'}(a', b') & = &
\det\nolimits_\kappa(b'|_{A'/(a')})^{-1}
\det\nolimits_\kappa(a'|_{A'/(b')})
\end{matrix}
$$
Hence we conclude that
$$
(-1)^{\text{length}_A(A'/(b'))}
d^A_{A'}(a', b' - a')
=
d^A_{A'}(b', b' - a') d^A_{A'}(a', b')
$$
the sign coming from the $-a'$ in the second equality above.
On the other hand, by Lemma \ref{lemma-periodic-determinant-sign} we have
$d^A_{A'}(b', b') = (-1)^{\text{length}_A(A'/(b'))}$ and the lemma
is proved.
\end{proof}

\end{document}
